La metodología del presente trabajo se basa en un enfoque de diseño conceptual y ilustración funcional mediante simulación, siguiendo etapas adaptadas del proceso de diseño de ingeniería. Estas etapas permiten estructurar el análisis de requerimientos, el modelado del sistema, y la evaluación de su funcionamiento a nivel virtual. Las fases metodológicas que se seguirán son las siguientes:

\section{Recolección y Análisis de Información}

La primera fase metodológica consiste en realizar una revisión exhaustiva de la información disponible sobre sistemas similares y las tecnologías involucradas en el desarrollo del kit educativo propuesto.

Esta etapa incluye:

\begin{itemize}
    \item Revisión documental de kits educativos similares, especialmente el myCobot AI Kit de Elephant Robotics, con el fin de identificar características funcionales, arquitecturas de hardware y software, así como elementos didácticos que puedan servir como referencia para el diseño propuesto.
    
    \item Estudio técnico del robot PhantomX Pincher, analizando sus especificaciones mecánicas, capacidades de movimiento, compatibilidad con servomotores Dynamixel, y limitaciones operativas que deban ser consideradas en el diseño del kit.
    
    \item Análisis de las capacidades de la Raspberry Pi 5 como plataforma de procesamiento embebido, evaluando su compatibilidad con ROS2, rendimiento computacional para tareas de visión por computador, y requisitos de integración con periféricos como cámaras y controladores de servomotores.
    
    \item Identificación de requerimientos funcionales y técnicos del sistema a diseñar, estableciendo las funcionalidades mínimas que debe cumplir el kit educativo y las restricciones técnicas derivadas de las plataformas hardware seleccionadas.
\end{itemize}

\section{Definición de Requerimientos}

Con base en la información recopilada en la fase anterior, se procede a establecer de manera formal los requerimientos del sistema. Esta fase es fundamental para orientar las decisiones de diseño posteriores.

Los aspectos a definir incluyen:

\begin{itemize}
    \item Especificación de los componentes de hardware necesarios, incluyendo cámaras, sistemas de sujeción (gripper o ventosa), fuentes de alimentación, y elementos de montaje mecánico.
    
    \item Especificación de los componentes de software requeridos, definiendo las bibliotecas de ROS2 necesarias, frameworks de visión por computador (como OpenCV), y herramientas de desarrollo y simulación.
    
    \item Delimitación de funcionalidades mínimas deseadas para el kit educativo, estableciendo qué tipos de objetos debe ser capaz de clasificar, el nivel de precisión esperado, y las operaciones básicas de manipulación que debe realizar.
    
    \item Identificación de restricciones físicas, considerando las dimensiones del espacio de trabajo, el peso máximo de los objetos a manipular, y las limitaciones cinemáticas del robot PhantomX Pincher.
    
    \item Identificación de restricciones computacionales, evaluando los recursos de procesamiento disponibles en la Raspberry Pi 5 y estableciendo compromisos entre velocidad de procesamiento y precisión en las tareas de visión.
    
    \item Identificación de restricciones pedagógicas, asegurando que el kit sea adecuado para su uso en entornos educativos, con una curva de aprendizaje apropiada y posibilidades de extensión para proyectos más avanzados.
\end{itemize}

\section{Diseño Conceptual del Sistema}

Esta fase constituye el núcleo del trabajo, donde se desarrolla la propuesta completa del kit educativo a nivel conceptual. El diseño se aborda desde dos perspectivas complementarias: hardware y software.

\subsection{Diseño de Hardware}

El diseño del hardware incluye:

\begin{itemize}
    \item Propuesta de arquitectura general del sistema, estableciendo la distribución espacial de componentes y sus interconexiones físicas.
    
    \item Esquema de conexión entre el robot PhantomX Pincher, la Raspberry Pi 5, la cámara, y los componentes de alimentación eléctrica.
    
    \item Diseño físico preliminar del kit, incluyendo el sistema de montaje de la cámara en posición óptima para la captura de imágenes del área de trabajo.
    
    \item Especificación del mecanismo de sujeción (gripper o ventosa) más adecuado según el tipo de objetos a manipular.
    
    \item Diseño de la disposición de componentes electrónicos, considerando aspectos de accesibilidad, seguridad, y facilidad de mantenimiento.
\end{itemize}

\subsection{Diseño de Software}

El diseño del software se centra en la arquitectura basada en ROS2:

\begin{itemize}
    \item Definición de la arquitectura de nodos ROS2 para la clasificación de objetos mediante visión por computador, estableciendo la estructura modular del sistema de software.
    
    \item Especificación del nodo de captura de imagen, responsable de la adquisición de imágenes desde la cámara y su publicación en topics de ROS2.
    
    \item Especificación del nodo de procesamiento visual, que implementa algoritmos de segmentación, detección de características, y preprocesamiento de imágenes.
    
    \item Especificación del nodo de clasificación de objetos, que aplica técnicas de reconocimiento de patrones o aprendizaje automático para identificar y categorizar los objetos detectados.
    
    \item Especificación del nodo de control del manipulador, que traduce las decisiones de clasificación en trayectorias de movimiento para el robot PhantomX Pincher.
    
    \item Representación esquemática de los flujos de datos entre percepción, procesamiento y acción, utilizando diagramas de arquitectura de ROS2 que muestren topics, servicios, y acciones.
\end{itemize}

\section{Modelado y Simulación del Sistema}

Para validar el diseño conceptual sin necesidad de construir el sistema físico completo, se emplea simulación en entornos compatibles con ROS2.

Esta fase incluye:

\begin{itemize}
    \item Modelado 3D preliminar de componentes del kit, utilizando herramientas CAD para representar la estructura física del sistema y verificar aspectos de ensamblaje y ergonomía.
    
    \item Configuración del entorno de simulación, empleando herramientas como Gazebo o RViz que son compatibles con ROS2, para crear un modelo virtual del espacio de trabajo.
    
    \item Implementación de nodos ROS2 básicos para la simulación, desarrollando versiones simplificadas de los nodos de captura de imagen, clasificación, y control.
    
    \item Simulación del flujo de operación mediante nodos de cámara virtual, clasificación básica basada en características simples de los objetos, y control del brazo robótico simulado.
    
    \item Observación del funcionamiento general del sistema en escenarios simulados, evaluando el comportamiento del sistema en diferentes condiciones operativas.
    
    \item Revisión del funcionamiento simulado, identificando posibles problemas de coordinación entre nodos, latencias en el flujo de datos, y oportunidades de optimización.
    
    \item Análisis de resultados de la simulación y ajustes al diseño, iterando sobre el diseño conceptual para incorporar las mejoras identificadas durante las pruebas de simulación.
\end{itemize}

\section{Construcción de Prototipo Físico Básico}

Como complemento al diseño conceptual y las simulaciones, se contempla la construcción de un prototipo físico básico que permita ilustrar de manera tangible los principios de funcionamiento del kit propuesto.

Las actividades de esta fase son:

\begin{itemize}
    \item Selección de componentes disponibles que permitan implementar una versión simplificada del diseño, priorizando la disponibilidad y el costo.
    
    \item Ensamblaje del prototipo físico básico, integrando el robot PhantomX Pincher con la Raspberry Pi 5, una cámara, y un sistema de sujeción simple.
    
    \item Implementación de software básico de control, desarrollando versiones funcionales pero simplificadas de los nodos de ROS2 diseñados.
    
    \item Pruebas de funcionamiento del prototipo, realizando experimentos con objetos reales para validar los conceptos de diseño.
    
    \item Documentación de las lecciones aprendidas durante la construcción y prueba del prototipo, que servirán como retroalimentación para refinar el diseño conceptual.
\end{itemize}

\section{Documentación y Recomendaciones}

La fase final de la metodología se centra en consolidar todo el trabajo realizado en documentación técnica de alta calidad que pueda servir como base para futuras implementaciones.

Esta fase comprende:

\begin{itemize}
    \item Elaboración de planos detallados del diseño mecánico, incluyendo vistas isométricas, cotas, y especificaciones de materiales.
    
    \item Creación de diagramas de bloques del sistema de hardware, mostrando las conexiones eléctricas y de comunicación entre componentes.
    
    \item Desarrollo de diagramas de arquitectura de software, documentando la estructura de nodos de ROS2, los flujos de mensajes, y las interfaces entre módulos.
    
    \item Redacción de guías para un posible uso académico del sistema en el futuro, incluyendo manuales de usuario, tutoriales de programación, y ejercicios prácticos sugeridos.
    
    \item Propuestas para adaptación pedagógica del kit en entornos educativos, considerando diferentes niveles de dificultad y objetivos de aprendizaje.
    
    \item Elaboración de estructura básica de guías didácticas, con actividades progresivas que permitan a los estudiantes familiarizarse con el sistema de manera gradual.
    
    \item Redacción de documentación técnica detallada, incluyendo especificaciones completas de hardware y software, procedimientos de calibración, y protocolos de mantenimiento.
    
    \item Identificación de posibles mejoras o extensiones futuras del diseño, proponiendo líneas de desarrollo que puedan ser exploradas en trabajos posteriores.
    
    \item Recomendaciones para la implementación práctica del kit, considerando aspectos logísticos, presupuestarios, y de integración curricular.
\end{itemize}

\section{Cronograma de Actividades}

El desarrollo del trabajo de grado se estructura en un período de 16 semanas, distribuidas según el siguiente cronograma:

\begin{table}[H]
\centering
\caption{Cronograma de actividades (Semanas 1-8)}
\label{tab:cronograma_1}
\footnotesize
\begin{tabular}{|p{6cm}|c|c|c|c|c|c|c|c|}
\hline
\textbf{Actividad} & \textbf{1} & \textbf{2} & \textbf{3} & \textbf{4} & \textbf{5} & \textbf{6} & \textbf{7} & \textbf{8} \\
\hline
Revisión bibliográfica y análisis del estado del arte & X & X & & & & & & \\
\hline
Identificación de requerimientos funcionales y técnicos & & & X & X & & & & \\
\hline
Diseño conceptual del hardware del kit & & & & & X & X & X & \\
\hline
Diseño de la arquitectura del sistema en ROS2 & & & & & & & X & X \\
\hline
Modelado 3D preliminar de componentes del kit & & & & & & & & X \\
\hline
Construcción de prototipo del kit & & & & & & & & \\
\hline
Simulación del sistema en entorno virtual & & & & & & & & \\
\hline
Análisis de resultados y ajustes al diseño & & & & & & & & \\
\hline
Redacción final del informe & & & & & & & & \\
\hline
\end{tabular}
\end{table}

\begin{table}[H]
\centering
\caption{Cronograma de actividades (Semanas 9-16)}
\label{tab:cronograma_2}
\footnotesize
\begin{tabular}{|p{6cm}|c|c|c|c|c|c|c|c|}
\hline
\textbf{Actividad} & \textbf{9} & \textbf{10} & \textbf{11} & \textbf{12} & \textbf{13} & \textbf{14} & \textbf{15} & \textbf{16} \\
\hline
Revisión bibliográfica y análisis del estado del arte & & & & & & & & \\
\hline
Identificación de requerimientos funcionales y técnicos & & & & & & & & \\
\hline
Diseño conceptual del hardware del kit & & & & & & & & \\
\hline
Diseño de la arquitectura del sistema en ROS2 & & & & & & & & \\
\hline
Modelado 3D preliminar de componentes del kit & X & & & & & & & \\
\hline
Construcción de prototipo del kit & X & X & & & & & & \\
\hline
Simulación del sistema en entorno virtual & & X & X & & & & & \\
\hline
Análisis de resultados y ajustes al diseño & & & X & X & X & & & \\
\hline
Redacción final del informe & & & & & X & X & X & \\
\hline
\end{tabular}
\end{table}